% GENERATED FILE. Edit canonical lessons, then run npm run generate:books.
% canonical-source-hash: fnv1a64:acbde4910f401e58
% canonical-lessons: GU-C275-vanchto-hato, GU-C275-banvu, GU-R275-the-scene

\chapter{What was going on}
\label{ch:gu-what-was-going-on}
\hlchaptermodality{275}

\begin{chapteropening}
\textbf{By the end of this chapter:} \emph{I can say what was going on and what used to happen (hũ vā̃cto hato), set a scene against an event with jyāre … tyāre, and ask what happened with banvũ (śũ banyũ?).}

Complete GU-R275-the-scene: set a scene, say what happened in it, and ask what happened.
\end{chapteropening}

\section[\texorpdfstring{\gu{હું} \gu{વાંચતો} \gu{હતો} (hũ vā̃cto hato)}{હું વાંચતો હતો (hũ vā̃cto hato)}]{\gu{હું} \gu{વાંચતો} \gu{હતો} (hũ vā̃cto hato) — was reading, used to read}

\label{lesson:GU-C275-vanchto-hato}



\begin{quote}
Say \textquotedblleft{}I don't read\textquotedblright{}, the A2 way. (\emph{hũ nathī vā̃cto}.) Keep its \textbf{-\gu{તો}}. Then say \textquotedblleft{}I was at home\textquotedblright{}. (\emph{hũ ghare hato}.)
\end{quote}

\begin{grammarlens}[title={Grammar Lens: -\gu{તો} with \gu{હતો}}]
\begin{quote}
\textbf{\gu{હું} \gu{પુસ્તક} \gu{વાંચતો} \gu{હતો}.}
\end{quote}

\begin{quote}
\emph{hũ pustak vā̃cto hato.} — \textquotedblleft{}I was reading a book.\textquotedblright{}
\end{quote}

The stem takes \textbf{-\gu{તો}}, then comes \textbf{\gu{હતો}}, and both follow the person:

\noindent
\begin{tabularx}{\linewidth}{@{}>{\raggedright\arraybackslash}X>{\raggedright\arraybackslash}X@{}}
\toprule
\textbf{who} & \textbf{was reading} \\
\midrule
a man & \textbf{\gu{વાંચતો} \gu{હતો}} (\emph{vā̃cto hato}) \\
a woman & \textbf{\gu{વાંચતી} \gu{હતી}} (\emph{vā̃ctī hatī}) \\
a neuter thing & \textbf{\gu{વાંચતું} \gu{હતું}} (\emph{vā̃ctũ hatũ}) \\
several, or respect & \textbf{\gu{વાંચતા} \gu{હતા}} (\emph{vā̃ctā hatā}) \\
\bottomrule
\end{tabularx}

There is no \textbf{-\gu{એ}} here, even with an object: \textbf{\gu{મીરા} \gu{ચા} \gu{પીતી} \gu{હતી}} (\emph{Mīrā cā pītī hatī}), \textquotedblleft{}Mira was drinking tea\textquotedblright{}. The doer form goes with the plain past, \textbf{\gu{પીધી}}.

The same form tells what used to happen: \textbf{\gu{પહેલાં} \gu{અરુણ} \gu{રોજ} \gu{ચાલતો} \gu{હતો}} (\emph{pahelā̃ Aruṇ roj cālto hato}), \textquotedblleft{}Arun used to walk every day\textquotedblright{}.
\end{grammarlens}

\subsection*{You'll want to know: the scene and the event}
\begin{quote}
\textbf{\gu{હું} \gu{વાંચતો} \gu{હતો} \gu{ત્યારે} \gu{ફોન} \gu{આવ્યો}.}
\end{quote}

\begin{quote}
\emph{hũ vā̃cto hato tyāre phon āvyo.} — \textquotedblleft{}I was reading when the phone rang.\textquotedblright{}
\end{quote}

\textbf{\gu{વાંચતો} \gu{હતો}} is the scene: it was going on. \textbf{\gu{આવ્યો}} is the event that broke in. \textbf{\gu{ત્યારે}} (\emph{tyāre}), \textquotedblleft{}at that time\textquotedblright{}, joins them, and \textbf{\gu{જ્યારે}} can open the scene: \textbf{\gu{જ્યારે} \gu{અરુણ} \gu{ઘરે} \gu{પહોંચ્યો}, \gu{ત્યારે} \gu{મીરા} \gu{રાંધતી} \gu{હતી}} (\emph{jyāre Aruṇ ghare pahõcyo, tyāre Mīrā rā̃dhtī hatī}), \textquotedblleft{}when Arun got home, Mira was cooking\textquotedblright{}.

\subsection*{Your turn}
Say these aloud:

\begin{itemize}
\raggedright
  \item \textquotedblleft{}It was raining\textquotedblright{} — \emph{varsād paḍto hato}
  \item \textquotedblleft{}Mira was cooking when Arun came home\textquotedblright{} — \emph{Mīrā rā̃dhtī hatī tyāre Aruṇ ghare āvyo}
\end{itemize}

\begin{tcolorbox}[breakable,colback=teal!4,colframe=teal!35!black,title={Before you move on}]
\textquotedblleft{}She was reading\textquotedblright{}? (\textbf{\gu{તે} \gu{વાંચતી} \gu{હતી}}.) Which part is the scene in \textbf{\gu{હું} \gu{વાંચતો} \gu{હતો} \gu{ત્યારે} \gu{ફોન} \gu{આવ્યો}}? (\textbf{\gu{વાંચતો} \gu{હતો}}.) \textquotedblleft{}Mira was drinking tea\textquotedblright{}: \textbf{\gu{મીરા}} or \textbf{\gu{મીરાએ}}? (\textbf{\gu{મીરા}}.)
\end{tcolorbox}
\section[\texorpdfstring{\gu{બનવું} (banvũ)}{બનવું (banvũ)}]{\gu{બનવું} (banvũ) — to happen; to become}

\label{lesson:GU-C275-banvu}



\begin{quote}
Before the new one: ask \textquotedblleft{}what happened?\textquotedblright{} the A2 way. (\textbf{\gu{શું} \gu{થયું}?}, \emph{śũ thayũ?}) Then say \textquotedblleft{}it was raining\textquotedblright{}. (\emph{varsād paḍto hato}.)
\end{quote}

\subsection*{You'll want to know: \gu{બનવું}}
\textbf{\gu{બનવું}} — \emph{banvũ} — \textquotedblleft{}to happen; to become\textquotedblright{}.

\begin{quote}
\textbf{\gu{ગઈકાલે} \gu{રાતે} \gu{શું} \gu{બન્યું}?}
\end{quote}

\begin{quote}
\emph{gaīkāle rāte śũ banyũ?} — \textquotedblleft{}What happened last night?\textquotedblright{}
\end{quote}

\textbf{\gu{શું} \gu{થયું}?} is the everyday question. \textbf{\gu{બનવું}} is for an event, a thing that takes place, and it suits a story. It takes no object, so its past follows its subject, as \textbf{\gu{ગયો}} does: \textbf{\gu{બન્યું}}, \textbf{\gu{બન્યો}}, \textbf{\gu{બની}}.

Its second sense is \textquotedblleft{}to become\textquotedblright{}:

\begin{quote}
\textbf{\gu{અરુણ} \gu{દાક્તર} \gu{બન્યો}.}
\end{quote}

\begin{quote}
\emph{Aruṇ dāktar banyo.} — \textquotedblleft{}Arun became a doctor.\textquotedblright{}
\end{quote}

\begin{cousinweb}
\textbf{\gu{બનવું}} (\emph{banvũ}) sits beside \textbf{\gu{બનાવવું}}, \textquotedblleft{}to make\textquotedblright{}, which you know: one makes a thing, the other is what a thing does when it comes about. CDIAL lists it under Sanskrit \emph{vanati}, \textquotedblleft{}desires, gains, makes ready\textquotedblright{}, and marks the Gujarati verb as taken from Hindi \emph{bannā} (CDIAL 11260).
\end{cousinweb}

\subsection*{Your turn}
Say these aloud:

\begin{itemize}
\raggedright
  \item \emph{banvũ}, then its past, \emph{banyũ}
  \item \textquotedblleft{}What happened last night?\textquotedblright{} — \emph{gaīkāle rāte śũ banyũ?}
  \item \textquotedblleft{}This used to happen every day\textquotedblright{} — \emph{āvũ roj bantũ hatũ}
\end{itemize}

\begin{tcolorbox}[breakable,colback=teal!4,colframe=teal!35!black,title={Before you move on}]
What does \textbf{\gu{બનવું}} mean? (\textquotedblleft{}To happen; to become\textquotedblright{}.) \textquotedblleft{}Arun became a doctor\textquotedblright{}? (\textbf{\gu{અરુણ} \gu{દાક્તર} \gu{બન્યો}}.) Which question is the everyday one? (\textbf{\gu{શું} \gu{થયું}?})
\end{tcolorbox}
\section[(review)]{Review — what was going on}

\label{lesson:GU-R275-the-scene}



\begin{quote}
Say \textquotedblleft{}Arun has gone home\textquotedblright{}. (\emph{Aruṇ ghare gayo che}.) Then \textquotedblleft{}I was going to go, but it rained\textquotedblright{}. (\emph{hũ javāno hato, paṇ varsād paḍyo}.)
\end{quote}

\subsection*{Your turn}
Say these aloud:

\begin{itemize}
\raggedright
  \item I was reading a book. Mira was drinking tea.
  \item I was reading when the phone rang.
  \item When Arun got home, Mira was cooking.
  \item Arun used to walk every day.
  \item What happened last night? Arun became a doctor.
\end{itemize}

\begin{tcolorbox}[breakable,colback=teal!4,colframe=teal!35!black,title={Before you move on}]
Which form paints the scene, \textbf{\gu{વાંચતો} \gu{હતો}} or \textbf{\gu{આવ્યો}}? (\textbf{\gu{વાંચતો} \gu{હતો}}.) Which word joins it to the event? (\textbf{\gu{ત્યારે}}.) \textquotedblleft{}What happened?\textquotedblright{}, in a story? (\textbf{\gu{શું} \gu{બન્યું}?})
\end{tcolorbox}
